如果集合~$A$ 有限，其大小可以用自然数来衡量；而比较集合的大小也就是比较这些自然数。如果集合是无穷的，情况就更复杂了。无穷的第一个层级是\emph{可数无穷}集合。集合~$A$ 是\emph{可数的}，若其元素可以排成一个\emph{枚举}——一个单向的无穷列表，即存在满射函数~$f\colon \PosInt \to A$。若集合可数但非有限，则它是\emph{可数无穷的}。Cantor的\emph{之字形法}表明，可数无穷集合的元素配对所成的集合也是可数的；由此可进一步证明，有理数集~$\Rat$ 也是可数的。

然而，也存在不可数的无穷集合：这些集合被称为\emph{不可数的}。证明一个集合不可数有两种方法：直接使用\emph{对角线论证}，或通过\emph{归约法}。要给出对角线论证，我们假设所讨论的集合~$A$ 是可数的，并利用一个假想的枚举来定义~$A$ 的一个元素，且由其定义方式可知，该元素必与枚举中的每个元素都不同。因此，该枚举终究未能穷尽~$A$ 的所有元素，由此便证明了~$A$ 的枚举不可能存在。归约法则是以满射的方式，将~$A$ 的每个元素与某个已知的不可数集合~$B$ 的元素相关联，从而证明~$A$ 是不可数的。若能做到这一点，那么由~$A$ 的一个假想枚举就能产生~$B$ 的一个枚举。由于~$B$ 不可数，故~$A$ 的枚举不可能存在。

一般而言，无穷集合可以按大小进行比较：若~$A$ 和~$B$ 之间存在双射，则它们大小相同，或者说\emph{等势}。我们也可以定义，若存在从~$A$ 到~$B$ 的单射函数，则~$A$ 不大于~$B$（$\cardle{A}{B}$）。根据 Schröder-Bernstein 定理，这实际上给出了无穷集合按大小排列的序。最后，Cantor 定理指出，对于任意~$A$，都有 $\cardless{A}{\Pow{A}}$。这推广了关于 $\Pow{\PosInt}$ 不可数的结果，并且表明无穷的层级不止两个，而是有无穷多个。